%% toc: P3 37 37.2,37.3,37.4
\chapter{Barwise's Axioms and the Karp Route to Lindstr\"om's Theorem}
\label{sec:p03-w03-p03-048-barwise}

Barwise's paper ``Axioms for abstract model theory'' proves Lindstr\"om's
characterization of first-order logic by a route of its own.  A logic with
the L\"owenheim property is shown to have the Karp property (partially
isomorphic structures satisfy the same sentences); a logic with the Karp
property that is stronger than first-order logic is shown to make the
ordinal \(\omega\) accessible; and an accessible \(\omega\) yields a
\(\Sigma^1_1\) definition of \((\omega,<)\), which is incompatible with
countable compactness and, given the L\"owenheim property, with the upward
L\"owenheim--Skolem property \cite[Part~II, \S\S1--3,
pp.~248--259]{Barwise1974Axioms}.  The shared proof of
\cref{thm:p03-lindstrom} in \cref{sec:p03-lindstrom} does not pass through
the Karp property or through accessible ordinals.

\emph{Purpose.} This chapter develops Barwise's Part~II notions with
individual guarded translations and restriction coherence. The full
composition requirement of the printed syntax is replaced by the
coherent interface of \cref{def:p03-w03-p03-048-b74-logic};
\cref{thm:p02-translation-composition} proves why this qualification is
necessary. Under this interface the chapter proves the conclusions of
his Theorems~1.2, 1.4 and 2.2, both directions
of his Theorem~3.2 for every \(\kappa=\beth_\kappa\geq\omega\) (for
\(\kappa>\omega\) the direction from \(\mathcal L_{\kappa\omega}\geq
\mathcal L^{\ast}\) to boundedness rests on a bound of Lopez-Escobar that
Barwise cites, and the same bound is cited here), his Corollary~3.3, and his
Corollary~3.4 (Lindstr\"om's theorem) with both of its alternative
hypotheses.  It then carries the Karp route into the regular relational
package of \cref{sec:p03-lindstrom} and derives \cref{thm:p03-lindstrom}
along it.  Where a construction is printed in the paper only in outline, or
where the printed construction needs an additional clause, the
reconstruction used here is stated and the difference is recorded.

\emph{Inputs.}  Barwise's languages, interpretations, pull-backs, categories
of languages, together with the guarded translation interface stated in
\cref{def:p03-w03-p03-048-b74-languages,def:p03-w03-p03-048-b74-logic};
first-order satisfaction, including on the empty structure; the compactness
theorem for first-order logic \cref{thm:fol-compactness} for languages of
any size.

\emph{Conclusions.}
\Cref{thm:p03-w03-p03-048-finite-omega} (Barwise's Theorem~1.2),
\cref{thm:p03-w03-p03-048-sup-accessible} (Theorem~1.4),
\cref{thm:p03-w03-p03-048-lowenheim-karp} (Theorem~2.2),
\cref{thm:p03-w03-p03-048-karp-omega} (Theorem~3.2 for \(\kappa=\omega\)),
\cref{thm:p03-w03-p03-048-lindstrom} (Corollary~3.4),
\cref{thm:p03-w03-p03-048-karp-kappa,thm:p03-w03-p03-048-kappa-forward}
(Theorem~3.2 for \(\kappa>\omega\) and Corollary~3.3), and, in the host
package, \cref{thm:p03-w03-p03-048-host-karp,thm:p03-w03-p03-048-host-route}.
The supporting lemmas are
\cref{thm:p03-w03-p03-048-basic,thm:p03-w03-p03-048-first-order,thm:p03-w03-p03-048-iso-closed,thm:p03-w03-p03-048-levels,thm:p03-w03-p03-048-fo-karp,thm:p03-w03-p03-048-countable,thm:p03-w03-p03-048-levels-transfinite}.

\emph{Source.}  The presentation follows the publisher's version of record,
\emph{Annals of Mathematical Logic} 7 (1974), pp.~221--265.  Barwise refers
the proofs of his preliminaries on partial isomorphisms to an earlier survey
\cite[Part~I, \S10, p.~243]{Barwise1974Axioms}; the facts of that kind used
below are proved here
(\cref{thm:p03-w03-p03-048-levels,thm:p03-w03-p03-048-fo-karp,thm:p03-w03-p03-048-countable}).

\section{Guarded translations}
\label{sec:p03-w03-p03-048-standing}

\begin{hbdefinition}{def:p03-w03-p03-048-b74-languages}{Languages and interpretations in Barwise's category}
Every language \(L\) is a set of symbols of finite arity containing
the distinguished unary relation symbol \(\mathbf V\). A constant
is a function symbol of arity zero. An \(L\)-structure interprets all
its symbols on a universe \(M=\mathbf V^{\mathfrak M}\); relations
are subsets of the appropriate powers of \(M\), and functions are
total. A partial structure may leave symbols undefined and functions
partially defined, but \(\mathbf V\) must be defined. Empty universes
are allowed whenever the symbols can be interpreted; a constant
precludes an empty total structure. A term of a language here is a
\emph{closed} term built from its constants and function symbols
\cite[Part~I, \S3, pp.~228--229]{Barwise1974Axioms}.

A \(k\)-term interpretation \(\alpha:L\to K\), for a finite
\(k\geq0\), specifies closed \(K\)-terms
\(t_1^\alpha,\ldots,t_k^\alpha\) and a map on symbols. It sends
an \(n\)-ary relation or function to an \((n+k)\)-ary symbol of
the same kind, and only \(\mathbf V\) may be sent to
\(\mathbf V\). A \(0\)-term interpretation fixing \(\mathbf V\)
is \emph{simple}; an injective simple interpretation is an
\emph{embedding}. Composition composes the symbol maps: if
\(\beta:K\to J\) has terms \(u_1,\ldots,u_\ell\), the composite
has the \(k+\ell\) terms
\((t_1^\alpha)^\beta,\ldots,(t_k^\alpha)^\beta,u_1,\ldots,u_\ell\).
Here term translation applies the symbol map and appends the parameters
\(u_1,\ldots,u_\ell\) at each function application. Identity maps
have no parameters. These operations define the category of all
languages and interpretations \cite[Part~I, \S4,
pp.~229--231]{Barwise1974Axioms}.

For a \(K\)-structure \(\mathfrak N\), put
\(a_j=(t_j^\alpha)^{\mathfrak N}\). The pull-back
\(\mathfrak N^{-\alpha}\) has universe
\[
 N_\alpha=\{b\in N:(b,a_1,\ldots,a_k)
                         \in(\mathbf V^\alpha)^{\mathfrak N}\}.
\]
For a relation symbol \(R\) its interpretation on that universe is
the set of tuples \(\bar b\) satisfying
\((\bar b,\bar a)\in(R^\alpha)^{\mathfrak N}\). For a function
symbol \(f\), evaluate \((f^\alpha)^{\mathfrak N}(\bar b,\bar a)\)
and retain the value only if it lies in \(N_\alpha\); otherwise the
pull-back function is undefined there. This includes the nullary case.
The structure \(\mathfrak N\) is \emph{\(\alpha\)-invertible}
exactly when its pull-back is a total \(L\)-structure. Inclusions
give reducts. Sending \(\mathbf V\) to a new unary predicate and
fixing the other symbols gives relativization to that predicate,
provided its interpretation contains the constants and is closed under
the functions \cite[Part~I, \S5, p.~231]{Barwise1974Axioms}.

A \emph{category of languages} \(\mathcal C\) is a subcategory of
this category satisfying the following five conditions
\cite[Part~I, \S6, Definition~6.2, pp.~233--234]{Barwise1974Axioms}:
\begin{enumerate}
\item It contains \(L\cap K\) and \(L\cup K\) whenever it contains
      \(L,K\).
\item Every inclusion between its languages is one of its morphisms.
\item Every finite language has an isomorphic copy in \(\mathcal C\),
      and, for each \(K\in\mathcal C\), such a copy whose only symbol
      in common with \(K\) is \(\mathbf V\).
\item Every \(0\)-term interpretation between its languages that fixes
      all but finitely many symbols is one of its morphisms.
\item For every \(L\in\mathcal C\), there are a constant \(d\notin L\),
      a \(K\in\mathcal C\) containing \(d\), and an injective
      \(1\)-term interpretation \(L\to K\) with term \(d\) that is
      a morphism of \(\mathcal C\).
\end{enumerate}
Barwise calls these category conditions tentative; they are the
conditions used in his paper, not a characterization of every possible
notion of abstract logic.
\end{hbdefinition}

\begin{hbdefinition}{def:p03-w03-p03-048-b74-logic}{The guarded translation interface}
A logic \(\mathcal L^\ast\) on \(\mathcal C\) supplies a class
\(L^\ast\) for each language and a one-step map
\(\alpha^\ast:L^\ast\to K^\ast\) for each interpretation.
The maps form a functor on the simple zero-term interpretations.
For an inclusion \(i:L_0\subseteq L\), they satisfy
\[
 (\alpha i)^\ast(\varphi)=\alpha^\ast(i^\ast(\varphi))
 \qquad(\varphi\in L_0^\ast).
\]
Full functoriality for domain-changing interpretations is not asserted:
\cref{thm:p02-translation-composition} gives the finite obstruction.
This is the same guarded interface as the first three clauses and
restriction condition of \cref{def:p02-w02c-p02-003-b74-logic}; its
additional characterization assumptions are stated separately below.
Write \(\varphi^\alpha=\alpha^\ast(\varphi)\). Satisfaction relates a
structure to a sentence in a language for that structure, with the
following semantic clauses from the source
\cite[Part~I, \S7, pp.~234--235]{Barwise1974Axioms}:
\begin{enumerate}
\item \emph{Occurrence.} Each sentence \(\varphi\) belongs to a least
      language \(L_\varphi\in\mathcal C\), ordered by inclusion.
      For every inclusion \(i:L_\varphi\subseteq K\), the map
      \(i^\ast\) is the inclusion of sentence classes.
\item \emph{Isomorphism.} Isomorphic structures satisfy the same
      sentences of their common language.
\item \emph{Translation.} For every morphism
      \(\alpha:L_\varphi\to K\) and \(K\)-structure
      \(\mathfrak N\),
      \[
       \mathfrak N\models\varphi^\alpha
       \quad\Longleftrightarrow\quad
       \mathfrak N\text{ is \(\alpha\)-invertible and }
       \mathfrak N^{-\alpha}\models\varphi.
      \]
\end{enumerate}
These three axioms do not by themselves assert closure under Boolean
operations, quantification, or finite occurrence. The additional
hypotheses used in the characterization theorems are stated separately
in \cref{def:p03-w03-p03-048-standing}.
\end{hbdefinition}

The source's syntax, semantics and domain-changing interpretations are
kept distinct in \cref{def:p03-w03-p03-048-b74-languages,def:p03-w03-p03-048-b74-logic}.
An individual relativization to a predicate \(V\) holds exactly when
that predicate is closed under the occurring functions, contains the
occurring constants, and the induced structure satisfies the sentence.
The explicit first-order construction is
\cref{thm:p02-guarded-fo-translation}. A second relativization retains its
own domain and closure conditions; the intermediate domain is not erased
by componentwise symbol composition. The universe symbol is
\(\mathbf V\), and empty structures are allowed
\cite[Part~I, \S3, p.~228]{Barwise1974Axioms}.

\begin{hbconvention}{def:p03-w03-p03-048-standing}{Standing hypotheses of Barwise's Part~II, \S\S1--3, and the readings used}
\begin{enumerate}
\item \emph{Closure} \cite[Part~II, \S1, p.~248]{Barwise1974Axioms}.
      \(\mathcal L^{\ast}\) contains all atomic sentences, is closed under
      conjunction and negation, and is closed under the existential
      quantifier; every language mentioned belongs to \(\mathcal C\).  The
      existential quantifier is an operation \(q\) defined on pairs
      \((c,\varphi)\) with \(c\) a constant symbol of \(L\) and
      \(\varphi\in L^{\ast}\); \(c\) does not occur in \(q(c,\varphi)\), and
      \(\mathfrak M\models q(c,\varphi)\) iff some \(a\) in the universe of
      \(\mathfrak M\) gives \((\mathfrak M,a)\models\varphi\) \cite[Part~I,
      \S9, 9.3, 9.7, pp.~239, 241]{Barwise1974Axioms}.  We read this truth
      condition for every structure \(\mathfrak M\) for a language of
      \(\mathcal C\) that contains the occurrence language of
      \(q(c,\varphi)\), with \((\mathfrak M,a)\) the structure that adds
      \(c\), if necessary, and interprets it by \(a\); on an empty
      structure \(q(c,\varphi)\) is then false.
\item \emph{Finite occurrence} \cite[Part~II, \S2, p.~254]{Barwise1974Axioms}.
      In \S\S2--3 every occurrence language \(L_\varphi\) is finite.
\item \emph{Partial isomorphism} \cite[Part~I, \S10, 10.1,
      pp.~243--244]{Barwise1974Axioms}.  A partial isomorphism from an
      \(L\)-structure \(\mathfrak M\) to an \(L\)-structure \(\mathfrak N\) is
      a one-one map \(F\) from a set \(D\subseteq M\) onto a set
      \(E\subseteq N\) that is an isomorphism of the partial substructures
      on \(D\) and \(E\), the partial functions being treated as their
      graphs: for each \(n\)-ary relation symbol \(R\) and
      \(\bar a\in D^{n}\), \(R^{\mathfrak M}(\bar a)\) iff
      \(R^{\mathfrak N}(F\bar a)\); for each \(n\)-ary function symbol \(f\)
      (\(n\geq0\)), \(\bar a\in D^{n}\) and \(b\in D\),
      \(f^{\mathfrak M}(\bar a)=b\) iff \(f^{\mathfrak N}(F\bar a)=F(b)\).
      \(I:\mathfrak M\cong_{p}\mathfrak N\) means that \(I\) is a set of
      partial isomorphisms with the back and forth property: for
      \(F\in I\) and \(a\in M\) (resp.\ \(b\in N\)) some \(G\in I\) has
      \(F\subseteq G\) and \(a\in\operatorname{dom}G\) (resp.\
      \(b\in\operatorname{ran}G\)).  We require \(I\) to be nonempty, as
      Barwise does for the sets \(I_\beta\) of his Theorem~10.5 (p.~245);
      with \(I=\emptyset\) the condition would hold between any two
      structures.  \(\mathfrak M\cong_{p}\mathfrak N\) means that some such
      \(I\) exists.
\item \emph{Equivalence and strength} \cite[Part~II, \S2, p.~254; \S3,
      3.1, p.~257]{Barwise1974Axioms}.  \(\mathfrak M\equiv^{\ast}\mathfrak N\)
      means that the \(L\)-structures \(\mathfrak M,\mathfrak N\) are models
      of the same sentences of \(L^{\ast}\).  For logics
      \(\mathcal L^{\ast},\mathcal L^{\#}\) on \(\mathcal C\),
      \(\mathcal L^{\#}\geq\mathcal L^{\ast}\) (\(\mathcal L^{\#}\) is as
      strong as \(\mathcal L^{\ast}\)) when every
      \(\mathcal L^{\ast}\)-sentence \(\varphi\) has an
      \(\mathcal L^{\#}\)-sentence \(\psi\) whose symbols all occur in
      \(\varphi\) and which has the same \(L_\varphi\)-structures as models;
      \(\mathcal L^{\ast}\equiv\mathcal L^{\#}\) when each is as strong as
      the other.  For \(\mathcal L_{\omega\omega}\geq\mathcal L^{\ast}\) the
      sentence \(\psi\) is a first-order sentence built from symbols of
      \(L_\varphi\), with \(=\) logical.
\end{enumerate}
\end{hbconvention}

\begin{hblemma}{thm:p03-w03-p03-048-basic}{proved=annex:p03-w03-p03-048-basic}{Reducts and invertible translations (Barwise, 7.2, 7.3, 7.5)}
Let \(\varphi\in L^{\ast}\) with \(L\) in \(\mathcal C\).
\begin{enumerate}
\item If \(L_\varphi\subseteq K\) and \(K\) is in \(\mathcal C\), then
      \(\varphi\in K^{\ast}\).
\item If \(\varphi\in K^{\ast}\cap K'^{\ast}\), \(\mathfrak M\) is a
      \(K\)-structure, \(\mathfrak M'\) is a \(K'\)-structure and
      \(\mathfrak M\restriction L_\varphi=\mathfrak M'\restriction L_\varphi\),
      then \(\mathfrak M\models\varphi\) iff \(\mathfrak M'\models\varphi\).
\item If \(\alpha:L\to K\) is a morphism of \(\mathcal C\) and the
      \(K\)-structure \(\mathfrak N\) is \(\alpha\)-invertible, then
      \(\mathfrak N\models\varphi^{\alpha}\) iff
      \(\mathfrak N^{-\alpha}\models\varphi\).
\end{enumerate}
\end{hblemma}

\begin{hblemma}{thm:p03-w03-p03-048-first-order}{proved=annex:p03-w03-p03-048-first-order}{First-order sentences inside \(\mathcal L^{\ast}\)}
Assume clause~1 of \cref{def:p03-w03-p03-048-standing}.  Let \(K\) be in
\(\mathcal C\), let \(\chi\) be a first-order formula of \(K\) whose free
variables are among the distinct variables \(x_1,\ldots,x_n\), and let
\(d_1,\ldots,d_n\) be distinct constant symbols not in \(K\) with
\(K(\bar d)\) in \(\mathcal C\).  There is \(\theta\in K(\bar d)^{\ast}\)
such that for every \(K\)-structure \(\mathfrak M\) (empty or not) and every
\(\bar a\in M^{n}\), \((\mathfrak M,\bar a)\models\theta\) iff
\(\mathfrak M\models\chi[\bar a]\).  If \(K\) is the set of symbols of
\(\chi\) together with \(\mathbf V\), every symbol occurring in \(\theta\) is
in \(K(\bar d)\).
\end{hblemma}

For a first-order sentence \(\chi\) of \(K\) we write \(\chi^{\bullet}\) for
the sentence of \(K^{\ast}\) given by the lemma with \(n=0\).  In
constructions below a first-order clause written as a conjunct of an
\(\mathcal L^{\ast}\)-sentence stands for its \(\chi^{\bullet}\), and a
clause with a free variable that is quantified over an
\(\mathcal L^{\ast}\)-sentence is formed with the operation \(q\) and a
constant symbol in place of the variable.  New symbols of prescribed kinds
and arities are taken, as clauses (iii) and (i) of the definition of a
category of languages allow, so that adjoining them to a language of
\(\mathcal C\) gives a language of \(\mathcal C\).

\section{\texorpdfstring{\(\Sigma^1_1\)}{Sigma-1-1} classes, accessible ordinals and the definability of \texorpdfstring{\((\omega,<)\)}{(omega,<)}}
\label{sec:p03-w03-p03-048-accessible}

\begin{hbdefinition}{def:p03-w03-p03-048-sigma}{\(\Sigma^1_1\) classes, accessible ordinals, boundedness (Barwise, Part~II, 1.1, 1.6)}
Let \(\mathcal K\) be a class of \(L\)-structures, \(L\) in \(\mathcal C\).
\begin{enumerate}
\item \(\mathcal K\) is \emph{\(\mathcal L^{\ast}\)-definable} if some
      \(\varphi\in L^{\ast}\) has \(\mathfrak M\in\mathcal K\) iff
      \(\mathfrak M\models\varphi\), for all \(L\)-structures \(\mathfrak M\).
      It is \emph{\(\Sigma^1_1\) in \(\mathcal L^{\ast}\)} if it is the class
      of pull-backs of an \(\mathcal L^{\ast}\)-definable class: there are a
      \(0\)-term morphism \(\gamma:L\to K\) and an
      \(\mathcal L^{\ast}\)-definable class \(\mathcal K'\) of
      \(K\)-structures, each \(\gamma\)-invertible, with
      \(\mathcal K=\{\mathfrak N^{-\gamma}:\mathfrak N\in\mathcal K'\}\)
      \cite[Part~II, \S1, 1.1, pp.~248--249]{Barwise1974Axioms}.
      Barwise requires \(\gamma\) to be the identity except possibly at
      \(\mathbf V\); we allow \(\gamma\) to be any one-one \(0\)-term
      morphism (a relativization combined with a renaming).  Barwise's own
      proofs of Theorems~1.2 and 1.4 use renamed order symbols (\(<'\) and
      \(\prec\), pp.~250--251), and the category axioms do not supply, for
      a given language, a copy in which a chosen symbol takes a prescribed
      name; the wider reading is the one those proofs establish.
      \(\mathcal K\) is \(\Pi^1_1\) if its complement among the
      \(L\)-structures is \(\Sigma^1_1\), and \(\Delta^1_1\) if it is both.
\item A structure \(\mathfrak M\) is definable (\(\Sigma^1_1\), etc.) if
      the class of structures isomorphic to it is.  Fix a binary relation
      symbol \(<\) with \(L_{<}=\{\mathbf V,<\}\) in \(\mathcal C\).  A class
      \(A\) of ordinals is \(\Sigma^1_1\) if the class
      \(\mathcal K_A=\{\mathfrak M:\mathfrak M\cong(\alpha,<)\text{ for some
      }\alpha\in A\}\) of \(L_{<}\)-structures is \(\Sigma^1_1\)
      \cite[1.1(ii)--(iii), p.~249]{Barwise1974Axioms}.
\item An ordinal \(\alpha\) is \emph{\(\mathcal L^{\ast}\)-accessible} if
      \(\alpha\in A\) for some \(\Sigma^1_1\) class \(A\) of ordinals
      \cite[p.~250]{Barwise1974Axioms}.  \(\mathcal L^{\ast}\) is
      \emph{bounded by \(\beta\)} if every \(\Sigma^1_1\) class of ordinals
      has all its members \(<\beta\); \(\operatorname{wo}(\mathcal L^{\ast})\)
      is the least such \(\beta\), when one exists \cite[p.~252,
      1.6]{Barwise1974Axioms}.  The strict reading of ``bounded by'' is the
      one under which Barwise's proof of Theorem~3.2 is valid: there an
      accessible \(\kappa\) must contradict
      \(\operatorname{wo}(\mathcal L^{\ast})\leq\kappa\).
\end{enumerate}
\end{hbdefinition}

\begin{hblemma}{thm:p03-w03-p03-048-iso-closed}{proved=annex:p03-w03-p03-048-iso-closed}{Classes of pull-backs are closed under isomorphism}
Let \(\gamma:L\to K\) be a one-one \(0\)-term morphism of \(\mathcal C\) and
let \(\psi\in K^{\ast}\) have only \(\gamma\)-invertible models.  Then
\(\{\mathfrak N^{-\gamma}:\mathfrak N\models\psi\}\) is closed under
isomorphism.  Consequently, if every pull-back of a model of \(\psi\) along a
morphism \(\gamma:L_{<}\to K\) is isomorphic to an ordinal, the class \(A\)
of the ordinals so obtained is \(\Sigma^1_1\).
\end{hblemma}

\begin{hbtheorem}{thm:p03-w03-p03-048-finite-omega}{proved=annex:p03-w03-p03-048-finite-omega}{Finite sets, \((\omega,<)\) and an infinite accessible ordinal (Barwise, Part~II, 1.2)}
Assume clause~1 of \cref{def:p03-w03-p03-048-standing}.  The following are
equivalent:
\begin{enumerate}
\item the class of finite \(\{\mathbf V\}\)-structures (finite sets, the
      empty one included) is \(\Sigma^1_1\) in \(\mathcal L^{\ast}\);
\item \((\omega,<)\) is \(\Sigma^1_1\) in \(\mathcal L^{\ast}\);
\item some \(\Sigma^1_1\) class of ordinals contains an infinite ordinal.
\end{enumerate}
\end{hbtheorem}

Barwise's statement and proof are at \cite[Part~II, \S1, 1.2,
pp.~249--250]{Barwise1974Axioms}.  His proof of
(iii)\(\Rightarrow\)(i) relativizes to a new unary \(U\) and requires \(U\)
to be an initial segment of the well-ordering in which every element with a
predecessor has an immediate predecessor.  An initial segment of type
\(\omega\) satisfies those clauses, so a model of type at least
\(\omega+1\) yields an infinite \(U\); the proof here adds the clause that
\(U\) is empty or has a largest element
(\cref{def:p03-w03-p03-048-finite-clause}).  Barwise notes that among the
logics of his Part~I, \S1, \((\omega,<)\) is definable in all but
\(\mathcal L_{\omega\omega}\) and \(\mathcal L(Q_1)\), in which it is not
\(\Sigma^1_1\) \cite[p.~250]{Barwise1974Axioms}.

\begin{hbtheorem}{thm:p03-w03-p03-048-sup-accessible}{proved=annex:p03-w03-p03-048-sup-accessible}{Downward closure and suprema of accessible ordinals (Barwise, Part~II, 1.4)}
Assume clause~1 of \cref{def:p03-w03-p03-048-standing}, and let \(A\) be a
\(\Sigma^1_1\) class of ordinals.
\begin{enumerate}
\item \(A'=\{\alpha:\alpha\leq\beta\text{ for some }\beta\in A\}\) is
      \(\Sigma^1_1\).
\item If \(A\) is a set and \(\beta=\sup A\), then \(\beta\) is
      \(\mathcal L^{\ast}\)-accessible.
\end{enumerate}
\end{hbtheorem}

By part~1, the \(\mathcal L^{\ast}\)-accessible ordinals form an initial
segment of the ordinals \cite[p.~251]{Barwise1974Axioms}.  Barwise prints
the sentence for part~2 with the clause that every predecessor of \(x\) is
below \(x\) in some indexed model; the proof here also requires that the
order on those predecessors agree with the order of that model, which is
what makes every proper initial segment of the new order a well-ordering.

\begin{hbwarning}{def:p03-w03-p03-048-finite-clause}{The extra clause in Theorem~1.2}
Let \(\psi\) define a class of well-orderings \((W,<)\) with a model of type
\(\omega+1\).  The three clauses printed on p.~250 (\(U\subseteq W\), \(U\)
closed downward, every element of \(U\) with a predecessor has an immediate
predecessor) hold when \(U\) is the initial segment of type \(\omega\).  So
the printed sentence has a model in which \(U\) is infinite, and the class of
its relativized reducts is not the class of finite sets.  With the added
clause ``\(U\) is empty or has a largest element'' every such \(U\) is a
finite ordinal, and every finite initial segment of an infinite model of
\(\psi\) is allowed.
\end{hbwarning}

\section{Partial isomorphism and the Karp property}
\label{sec:p03-w03-p03-048-karp}

\begin{hbdefinition}{def:p03-w03-p03-048-karp}{Karp and L\"owenheim properties (Barwise, Part~II, 2.1)}
\(\mathcal L^{\ast}\) has the \emph{Karp property} if for every \(L\) in
\(\mathcal C\) and all \(L\)-structures \(\mathfrak M,\mathfrak N\),
\(\mathfrak M\cong_{p}\mathfrak N\) implies
\(\mathfrak M\equiv^{\ast}\mathfrak N\) \cite[Part~II, \S2, 2.1,
p.~254]{Barwise1974Axioms}.  \(\mathcal L^{\ast}\) has the
\emph{L\"owenheim property} if every \(\mathcal L^{\ast}\)-sentence with a
model has a model of power at most \(\aleph_0\) \cite[p.~255]{Barwise1974Axioms}.
\end{hbdefinition}

For a finite language \(L\) and \(L\)-structures \(\mathfrak M,\mathfrak N\),
let \(J_0(\mathfrak M,\mathfrak N)\) be the set of finite partial
isomorphisms from \(\mathfrak M\) to \(\mathfrak N\), and let
\(J_{m+1}(\mathfrak M,\mathfrak N)\) be the set of \(p\in J_m\) such that for
each \(a\in M\) some \(q\in J_m\) with \(p\subseteq q\) has
\(a\in\operatorname{dom}q\), and for each \(b\in N\) some \(q\in J_m\) with
\(p\subseteq q\) has \(b\in\operatorname{ran}q\).  These are the finite
levels of Barwise's sequences \(I_0\supseteq I_1\supseteq\cdots\) of
Theorem~10.5 \cite[p.~245]{Barwise1974Axioms}.

\begin{hblemma}{thm:p03-w03-p03-048-levels}{proved=annex:p03-w03-p03-048-levels}{Finite-level types}
Let \(L\) be finite.  Each \(J_m(\mathfrak M,\mathfrak N)\) is closed under
subsets and \(J_{m+1}\subseteq J_m\).  For all \(m,n\) there is a finite set
\(\Phi_{m,n}\) of first-order formulas of \(L\) with free variables among
\(x_1,\ldots,x_n\), and for each \(L\)-structure \(\mathfrak M\) and
\(\bar a\in M^{n}\) a formula \(\tau^{m}_{\mathfrak M,\bar a}\in\Phi_{m,n}\),
such that \(\mathfrak M\models\tau^{m}_{\mathfrak M,\bar a}[\bar a]\) and, for
every \(\mathfrak N\) and \(\bar b\in N^{n}\),
\(\mathfrak N\models\tau^{m}_{\mathfrak M,\bar a}[\bar b]\) iff
\(\{(a_i,b_i):i\leq n\}\) is a map that belongs to
\(J_m(\mathfrak M,\mathfrak N)\).  Consequently, if \(\varphi\in L^{\ast}\),
\(L=L_\varphi\), and for some \(m\) there are no \(L\)-structures
\(\mathfrak M\models\varphi\), \(\mathfrak N\not\models\varphi\) with
\(\emptyset\in J_m(\mathfrak M,\mathfrak N)\), then \(\varphi\) has the same
\(L\)-structures as models as the first-order sentence
\(\bigvee\{\tau^{m}_{\mathfrak M}:\mathfrak M\models\varphi\}\) of \(L\).
\end{hblemma}

Barwise's own sentences \(\sigma^{\alpha}_{(\mathfrak M,a)}\) (p.~246) are
built from atomic and negated atomic sentences of quantifier rank \(0\), and
by his Definition~10.4 an atomic sentence in which a function symbol of
positive arity occurs has positive rank (p.~245).  For \(L=\{\mathbf V,f\}\)
with \(f\) unary, every \(\sigma^{\alpha}_{\mathfrak M}\) is then a sentence
about equality alone, while \(\forall x\,f(x)=x\), of rank \(2\) in his
count, separates \((\omega,\mathrm{id})\) from \((\omega,\mathrm{succ})\).
The formulas \(\tau^{m}\) above therefore also record the graphs of the
functions on the listed elements, which is what partial isomorphism in
clause~3 of \cref{def:p03-w03-p03-048-standing} preserves.

\begin{hblemma}{thm:p03-w03-p03-048-fo-karp}{proved=annex:p03-w03-p03-048-fo-karp}{First-order and infinitary logic have the Karp property}
If \(I:\mathfrak M\cong_{p}\mathfrak N\), then for every formula
\(\chi(x_1,\ldots,x_n)\) of \(\mathcal L_{\infty\omega}\)
(\cref{def:p03-w03-p03-048-infinitary}), in particular for every first-order
formula, every \(F\in I\) and every
\(\bar a\in(\operatorname{dom}F)^{n}\),
\(\mathfrak M\models\chi[\bar a]\) iff \(\mathfrak N\models\chi[F\bar a]\).
In particular \(\mathfrak M\) and \(\mathfrak N\) satisfy the same
sentences of \(\mathcal L_{\infty\omega}\).
\end{hblemma}

\begin{hblemma}{thm:p03-w03-p03-048-countable}{proved=annex:p03-w03-p03-048-countable}{Countable partially isomorphic structures are isomorphic (Barwise, Part~I, 10.2)}
Let \(\mathfrak M,\mathfrak N\) be countable \(L\)-structures and
\(I:\mathfrak M\cong_{p}\mathfrak N\).  For every \(F_0\in I\) there is an
isomorphism \(F:\mathfrak M\cong\mathfrak N\) with \(F_0\subseteq F\).
\end{hblemma}

Barwise states this also for countably generated structures
\cite[Part~I, \S10, 10.2, p.~244]{Barwise1974Axioms}; only the countable
case is used.  \Cref{thm:p03-w03-p03-039-countable-back-forth,thm:p03-w03-p03-039-karp}
state the countable back and forth and Karp's theorem for partial embeddings
that preserve every atomic formula, nested terms included
(\cref{def:p03-w03-p03-039-partial-iso}).  Barwise's partial isomorphisms
preserve only the graphs of the functions on their domains, and his
structures may be empty; the proofs here are written for his notion.

\begin{hbtheorem}{thm:p03-w03-p03-048-lowenheim-karp}{proved=annex:p03-w03-p03-048-lowenheim-karp}{The L\"owenheim property gives the Karp property (Barwise, Part~II, 2.2)}
Assume clauses~1 and 2 of \cref{def:p03-w03-p03-048-standing}.  If
\(\mathcal L^{\ast}\) has the L\"owenheim property, it has the Karp
property.
\end{hbtheorem}

\begin{hbwarning}{def:p03-w03-p03-048-karp-sentence}{The sentence of Theorem~2.2 as printed and as used}
Barwise's sentence (p.~255) is the conjunction of \(\varphi^{(U)}\),
\((\neg\varphi)^{(W)}\) and clauses (1)--(5) about a pairing function
\(p\) and a binary \(E\) on codes of tuples: \(p\) is injective,
\(E\)-related codes extend forth into \(W\) and back into \(U\), and
\(E\)-related codes of length \(n\) agree on each \(n\)-ary relation symbol
and on each graph of an \(n\)-ary function symbol.  Read literally, these
clauses do not make the map coded by a long tuple preserve the relations
among its shorter subtuples, do not make it preserve equality, and do not
assert that any code is \(E\)-related; and since both copies use the same
symbols, a constant symbol would have to name an element of \(U\) and of
\(W\) at once.  The proof here keeps the plan (relativize \(\varphi\) and
\(\neg\varphi\) to two copies, describe a back and forth system by a
sentence, apply the L\"owenheim property, then
\cref{thm:p03-w03-p03-048-countable}) and changes the coding: the second
copy uses renamed symbols, and the system is a family of maps indexed by
positions, with clauses quantifying over all pairs of a position, so
finitely many clauses suffice.  Barwise's remark that the converse fails
(\(\mathcal L_{\infty\omega}\) has the Karp property but not the
L\"owenheim property) and his Theorem~2.3 (the converse under
interpolation) are on p.~256 and are not used below.
\end{hbwarning}

\section{Characterizing logics with the Karp property}
\label{sec:p03-w03-p03-048-characterizing}

Barwise states Theorem~3.2 for every cardinal \(\kappa=\beth_\kappa\geq\omega\)
in his indexing of the \(\beth\) function
(\cref{def:p03-w03-p03-048-infinitary}): \(\mathcal L_{\kappa\omega}\geq
\mathcal L^{\ast}\) iff \(\mathcal L^{\ast}\) has the Karp property and
\(\operatorname{wo}(\mathcal L^{\ast})\leq\kappa\) \cite[Part~II, \S3, 3.2,
p.~257]{Barwise1974Axioms}.  The case \(\kappa=\omega\), which carries
Corollary~3.4, is \cref{thm:p03-w03-p03-048-karp-omega}; the cases
\(\kappa>\omega\) and Corollary~3.3 are
\cref{thm:p03-w03-p03-048-karp-kappa,thm:p03-w03-p03-048-kappa-forward} in
\cref{sec:p03-w03-p03-048-infinitary}.

\begin{hbtheorem}{thm:p03-w03-p03-048-karp-omega}{proved=annex:p03-w03-p03-048-karp-omega}{Karp property and boundedness by \(\omega\) (Barwise, Part~II, 3.2 for \(\kappa=\omega\))}
Assume clauses~1 and 2 of \cref{def:p03-w03-p03-048-standing}.  Then
\(\mathcal L_{\omega\omega}\geq\mathcal L^{\ast}\) iff \(\mathcal L^{\ast}\)
has the Karp property and is bounded by \(\omega\).
\end{hbtheorem}

\begin{hbcorollary}{thm:p03-w03-p03-048-lindstrom}{proved=annex:p03-w03-p03-048-lindstrom}{Guarded characterization}
Assume clauses~1 and 2 of \cref{def:p03-w03-p03-048-standing}.  If
\(\mathcal L^{\ast}\) has the L\"owenheim property and either
\begin{enumerate}
\item[(a)] (upward L\"owenheim--Skolem) every \(\mathcal L^{\ast}\)-sentence
      with an infinite model has an uncountable model, or
\item[(b)] (countable compactness) whenever \(T\) is a countable subset of
      \(L^{\ast}\) for some \(L\) in \(\mathcal C\) and every finite subset
      of \(T\) has a model, \(T\) has a model,
\end{enumerate}
then \(\mathcal L_{\omega\omega}\geq\mathcal L^{\ast}\): every
\(\mathcal L^{\ast}\)-sentence \(\varphi\) has the same
\(L_\varphi\)-structures as models, including every empty structure permitted by the signature,
as a first-order sentence built from symbols of \(L_\varphi\).  Conversely,
every first-order sentence \(\chi\) of a language \(L\) in \(\mathcal C\) has
the same \(L\)-structures as models as \(\chi^{\bullet}\in L^{\ast}\), whose
symbols lie in \(L_\chi\cup\{\mathbf V\}\) whenever that set is a language of
\(\mathcal C\); in that case \(\mathcal L^{\ast}\equiv\mathcal L_{\omega\omega}\).
\end{hbcorollary}

The route of the proof is Barwise's: Theorem~2.2 gives the Karp property;
Theorem~3.2 makes \(\omega\) accessible if \(\mathcal L^{\ast}\) is stronger
than first-order logic; Theorem~1.2 then makes \((\omega,<)\) a
\(\Sigma^1_1\) structure, and a \(\Sigma^1_1\) definition of \((\omega,<)\)
contradicts (b) with one new constant, and contradicts (a) once the
L\"owenheim property supplies a countable model and a new injection into the
ordered part \cite[Part~II, \S3, proof of 3.4, p.~259]{Barwise1974Axioms}.

\begin{hbwarning}{def:p03-w03-p03-048-versus-lindstrom}{Barwise's notion of logic and Lindstr\"om's}
Barwise records that his treatment removes the exceptions that Lindstr\"om's
formulation had to make for finite models, and he locates the difference
between the two notions of logic in two places
\cite[Part~II, \S3, p.~259]{Barwise1974Axioms}.  In one sense the only
difference is that his logics permit relativization.  More fundamentally,
Lindstr\"om works with classes of structures directly, whereas Barwise takes
the relation between sentences and the structures they define as the object
of study, so that closure conditions such as the formation of indexed unions
and their inverses appear as consequences of syntactic operations (the
translation axiom applied to \(1\)-term interpretations).  The proof of
\cref{thm:p03-w03-p03-048-lindstrom} uses exactly such an operation: the
indexed union in \cref{thm:p03-w03-p03-048-finite-omega} and
\cref{thm:p03-w03-p03-048-sup-accessible}.  Barwise also states that all
results of his Part~II are proved by monomorphisms \cite[Part~I, \S5,
note~3, p.~232]{Barwise1974Axioms}, and, where his Theorem~2.2 needs one
clause for each symbol of \(L_\varphi\), he refers to his problem~P5
\cite[Part~II, \S2, p.~255]{Barwise1974Axioms}.
\end{hbwarning}

\section{Theorem~3.2 for uncountable \texorpdfstring{\(\kappa\)}{kappa} and Corollary~3.3}
\label{sec:p03-w03-p03-048-infinitary}

\begin{hbdefinition}{def:p03-w03-p03-048-infinitary}{Barwise's \(\beth\) indexing and the logics \(\mathcal L_{\kappa\omega}\) (Part~I, \S10)}
\begin{enumerate}
\item For a cardinal \(\lambda\): \(\beth_0(\lambda)=\lambda\),
      \(\beth_{\alpha+1}(\lambda)=2^{\beth_\alpha(\lambda)}\), and
      \(\beth_\beta(\lambda)=\sup_{\alpha<\beta}\beth_\alpha(\lambda)\) for a
      limit \(\beta\); \(\beth_\alpha=\beth_\alpha(0)\)
      \cite[Part~I, \S10, p.~246]{Barwise1974Axioms}.  In this indexing
      \(\beth_\omega=\aleph_0=\omega\), so \(\kappa=\omega\) satisfies
      \(\kappa=\beth_\kappa\) (the more common indexing starts with
      \(\beth_0=\aleph_0\)).
\item \(\mathcal L_{\infty\omega}\) is first-order logic with conjunctions
      \(\bigwedge\Phi\) and disjunctions \(\bigvee\Phi\) of arbitrary sets
      \(\Phi\) of formulas; truth is Tarski's, \(\bigwedge\Phi\) holding when
      every member holds.  Only formulas whose free variables lie among
      finitely many variables are used below.  The set of subformulas is
      \(\operatorname{sub}(\varphi)=\{\varphi\}\) for atomic \(\varphi\),
      \(\{\varphi\}\cup\operatorname{sub}(\psi)\) for \(\varphi\) one of
      \(\forall x\,\psi\), \(\exists x\,\psi\), \(\neg\psi\), and
      \(\{\varphi\}\cup\bigcup_{\psi\in\Phi}\operatorname{sub}(\psi)\) for
      \(\bigwedge\Phi\), \(\bigvee\Phi\); \(|\varphi|\) is the cardinality of
      \(\operatorname{sub}(\varphi)\).  For a cardinal \(\kappa\geq\omega\),
      \(\mathcal L_{\kappa\omega}\) consists of the formulas with
      \(|\varphi|<\kappa\) \cite[p.~246]{Barwise1974Axioms}; so
      \(\mathcal L_{\omega\omega}\) is first-order logic.
\item \(\mathcal L_{\kappa\omega}\geq\mathcal L^{\ast}\) and
      \(\mathcal L_{\infty\omega}\geq\mathcal L^{\ast}\) are read as in
      clause~4 of \cref{def:p03-w03-p03-048-standing}, with
      \(\mathcal L_{\kappa\omega}\)- or \(\mathcal L_{\infty\omega}\)-sentences
      \(\psi\) built from the symbols of \(L_\varphi\).  A class of ordinals
      is \(\Sigma^1_1\) in \(\mathcal L_{\lambda\omega}\) when
      \cref{def:p03-w03-p03-048-sigma} holds with an
      \(\mathcal L_{\lambda\omega}\)-sentence in place of the
      \(\mathcal L^{\ast}\)-sentence.  \(\mathcal L^{\ast}\) is
      \emph{bounded} if every \(\Sigma^1_1\) class of ordinals is bounded,
      that is, is a set \cite[Part~II, \S1, p.~252]{Barwise1974Axioms}.
\end{enumerate}
The axiom of choice is assumed: cardinal arithmetic, the regularity of
successor cardinals and the comparison of cardinalities below use it.
\end{hbdefinition}

For a finite language \(L\) the sets \(J_m\) of
\cref{sec:p03-w03-p03-048-karp} are continued through the ordinals:
\(J_0(\mathfrak M,\mathfrak N)\) is the set of finite partial isomorphisms,
\(J_{\beta+1}\) is obtained from \(J_\beta\) by the back and forth clause used
for \(J_{m+1}\), and \(J_\lambda=\bigcap_{\beta<\lambda}J_\beta\) for a limit
\(\lambda\).  These are the sequences \(I_\beta\) of Barwise's Theorem~10.5,
taken for finite partial isomorphisms \cite[p.~245]{Barwise1974Axioms}.

\begin{hblemma}{thm:p03-w03-p03-048-levels-transfinite}{proved=annex:p03-w03-p03-048-levels-transfinite}{Transfinite types and their count (Barwise, 10.6--10.9)}
Let \(L\) be finite.
\begin{enumerate}
\item Each \(J_\beta(\mathfrak M,\mathfrak N)\) is closed under subsets, and
      \(J_\gamma\subseteq J_\beta\) for \(\beta\leq\gamma\).
\item For every ordinal \(\beta\) and every \(n\) there is a set
      \(\Phi_{\beta,n}\) of \(\mathcal L_{\infty\omega}\)-formulas of \(L\)
      with free variables among \(x_1,\ldots,x_n\), and for each
      \(L\)-structure \(\mathfrak M\) and \(\bar a\in M^{n}\) a formula
      \(\tau^{\beta}_{\mathfrak M,\bar a}\in\Phi_{\beta,n}\), such that
      \(\mathfrak N\models\tau^{\beta}_{\mathfrak M,\bar a}[\bar b]\) iff
      \(\{(a_i,b_i):i\leq n\}\) is a map that belongs to
      \(J_\beta(\mathfrak M,\mathfrak N)\); for finite \(\beta\) these are the
      formulas of \cref{thm:p03-w03-p03-048-levels}.
\item \(|\Phi_{\beta,n}|\leq\beth_{\beta+1}(\aleph_0)\) and
      \(|\tau^{\beta}_{\mathfrak M,\bar a}|\leq\beth_{\beta+1}(\aleph_0)\).
      If \(\kappa=\beth_\kappa>\omega\) and \(\beta<\kappa\), both are
      \(<\kappa\).
\item If \(\varphi\in L^{\ast}\), \(L=L_\varphi\), and for some \(\beta\)
      there are no \(L\)-structures \(\mathfrak M\models\varphi\),
      \(\mathfrak N\not\models\varphi\) with \(\emptyset\in
      J_\beta(\mathfrak M,\mathfrak N)\), then \(\varphi\) has the same
      \(L\)-structures as models as the \(\mathcal L_{\infty\omega}\)-sentence
      \(\bigvee\{\tau^{\beta}_{\mathfrak M}:\mathfrak M\models\varphi\}\) of
      \(L\), which lies in \(\mathcal L_{\kappa\omega}\) when
      \(\kappa=\beth_\kappa>\omega\) and \(\beta<\kappa\).
\end{enumerate}
\end{hblemma}

Barwise's sentences \(\sigma^{\alpha}_{(\mathfrak M,a)}\) and his count
(Lemma~10.8: fewer than \(\kappa\) of them for \(\alpha<\kappa\), each in
\(\mathcal L_{\kappa\omega}\)) are at \cite[pp.~246--247]{Barwise1974Axioms};
his proof of 10.8 for \(\kappa>\omega\) states the bound
\(\beth_{\alpha+1}(\lambda)\), with \(\lambda=|L|+\aleph_0\), and leaves the
induction to the reader.  The formulas \(\tau^{\beta}\) replace his
\(\sigma^{\alpha}\) for the reason given after
\cref{thm:p03-w03-p03-048-levels}.

\subsection{Ranks of small pointed models}
\label{sec:p03-infinitary-boundedness}

The fragment lemmas give small elementary submodels and elementary
unions. When a sentence permits only well-ordered interpretations,
these two facts turn elementary extensions of pointed models into
a well-founded relation on a set. Its rank gives a uniform ordinal
bound.

\begin{hbtheorem}{thm:p03-ib-rank}{proved=annex:p03.ib.1}{A bound from small pointed models}
Let \(\lambda\) be an infinite cardinal and let \(F\) be a fragment
in the sense of \cref{def:p03-w03-p03-039-fragment}, with
\(|F|\leq\lambda\) and vocabulary of size at most \(\lambda\).
Every formula has finitely many free variables. Suppose
\(\psi\in F\) is a sentence and \(\delta(x),\rho(x,y)\in F\),
and that in every model \(\mathfrak M\models\psi\), the restriction
of \(\rho^{\mathfrak M}\) well-orders
\(D^{\mathfrak M}=\{a:\mathfrak M\models\delta(a)\}\).
Then there is an ordinal \(\beta<(2^\lambda)^+\) such that all
these interpreted well-orders have order type less than \(\beta\).
Empty carriers are allowed and have type zero. The nonempty-model
fragment lemmas are applied only to models with a marked carrier
point.
\end{hbtheorem}

\begin{hbcorollary}{thm:p03-ib-individual}{proved=annex:p03.ib.2}{Individual bounds}
Every \(L_{\infty\omega}\)-sentence that makes a designated
binary relation a well-order of its field in all models has a
uniform ordinal bound on those order types. Nevertheless, for each
ordinal \(\alpha\geq2\), a sentence of \(L_{\infty\omega}\)
pins down \(\alpha\). Thus this logic is bounded in the
sentence-by-sentence sense while its well-ordering number is
\(w(L_{\infty\omega})=\infty\), as in
\cite[Ch.~II, Example~5.2.3(e), p.~62]{Ebbinghaus1985ExtendedLogics}.
\end{hbcorollary}

The next theorem applies the rank bound to the particular
renaming-and-relativization pullbacks of Barwise's framework.
Only the symbols occurring in the witnessing sentence and the
images of the carrier and order symbols are needed.


\begin{hbtheorem}{thm:p03-w03-p03-048-lopez-escobar}{proved=annex:p03.ib.3}{The Lopez-Escobar bound}
For every infinite regular cardinal \(\lambda\), every class of ordinals
that is \(\Sigma^1_1\) in \(\mathcal L_{\lambda\omega}\) has all its
members below \((2^{\lambda})^{+}\); that is,
\(\operatorname{wo}(\mathcal L_{\lambda\omega})\leq(2^{\lambda})^{+}\).
\end{hbtheorem}

Barwise attributes this bound to Lopez-Escobar and refers to a paper of
Barwise and Kunen for a direct proof
\cite[Part~II, \S3, proof of 3.2, p.~257]{Barwise1974Axioms}.
The proof here derives it from the rank of small pointed models in
\cref{thm:p03-ib-rank}. It supplies the boundedness step in the forward
direction of \cref{thm:p03-w03-p03-048-kappa-forward}.

\begin{hbtheorem}{thm:p03-w03-p03-048-karp-kappa}{proved=annex:p03-w03-p03-048-karp-kappa}{From the Karp property to \(\mathcal L_{\kappa\omega}\) and \(\mathcal L_{\infty\omega}\) (Barwise, 3.2 and 3.3, \((\Leftarrow)\))}
Assume clauses~1 and 2 of \cref{def:p03-w03-p03-048-standing}, and that
\(\mathcal L^{\ast}\) has the Karp property.
\begin{enumerate}
\item If \(\kappa=\beth_\kappa\geq\omega\) and \(\mathcal L^{\ast}\) is
      bounded by \(\kappa\), then \(\mathcal L_{\kappa\omega}\geq
      \mathcal L^{\ast}\).
\item If \(\mathcal L^{\ast}\) is bounded, then
      \(\mathcal L_{\infty\omega}\geq\mathcal L^{\ast}\).
\end{enumerate}
\end{hbtheorem}

\begin{hbtheorem}{thm:p03-w03-p03-048-kappa-forward}{proved=annex:p03-w03-p03-048-kappa-forward}{From \(\mathcal L_{\kappa\omega}\) and \(\mathcal L_{\infty\omega}\) to the Karp property and boundedness (Barwise, 3.2 and 3.3, \((\Rightarrow)\))}
Let \(\mathcal L^{\ast}\) be a logic on \(\mathcal C\).
\begin{enumerate}
\item If \(\kappa=\beth_\kappa>\omega\) and \(\mathcal L_{\kappa\omega}\geq
      \mathcal L^{\ast}\), then \(\mathcal L^{\ast}\) has the Karp property
      and is bounded by \(\kappa\).
\item If \(\mathcal L_{\infty\omega}\geq\mathcal L^{\ast}\), then
      \(\mathcal L^{\ast}\) has the Karp property and is bounded.
\end{enumerate}
\end{hbtheorem}

With \cref{thm:p03-w03-p03-048-karp-omega} for \(\kappa=\omega\), the two
theorems give Barwise's Theorem~3.2 for every \(\kappa=\beth_\kappa\geq
\omega\), and part~2 of each gives his Corollary~3.3: \(\mathcal
L_{\infty\omega}\geq\mathcal L^{\ast}\) iff \(\mathcal L^{\ast}\) is bounded
and has the Karp property \cite[3.3, p.~258]{Barwise1974Axioms}.  Barwise
describes the corollary as a strengthening of a result of his earlier paper
on absolute logics.  The part of Theorem~3.2 that rests on
\cref{thm:p03-w03-p03-048-lopez-escobar} is the direction \(\Rightarrow\) for
\(\kappa>\omega\); the direction used for Lindstr\"om's theorem is proved
outright.

\section{Barwise's route and the shared characterization}
\label{sec:p03-w03-p03-048-bridge}

\begin{hbconvention}{def:p03-w03-p03-048-bridge}{Correspondence with \cref{thm:p03-lindstrom}}
\begin{enumerate}
\item \emph{Logics.}  The regular relational package of
      \cref{sec:p03-lindstrom} is stated for relational signatures and
      nonempty structures, with renaming and relativization to a fresh
      nonempty unary predicate.  Barwise's logics live on a category of
      languages with function and constant symbols, allow the empty
      structure, and obey the translation axiom for every morphism of the
      category, including \(1\)-term interpretations (indexed unions).
      Neither notion is a special case of the other.
\item \emph{Hypotheses.}  The host hypothesis~2 of \cref{thm:p03-lindstrom}
      is the L\"owenheim property; host hypothesis~1 is countable
      compactness for countable sets of sentences, where Barwise's (b)
      requires the set to lie in one \(L^{\ast}\).  Barwise also allows
      hypothesis~(a), upward L\"owenheim--Skolem, in place of compactness.
      Finite occurrence is assumed in both.
\item \emph{Proof routes.}  The shared proof in
      \cref{sec:p03-lindstrom} compares structures through finite partial
      isomorphisms and potential isomorphism
      (\cref{def:p03-potential-isomorphism}) and reaches a contradiction with
      compactness by a nonstandard index.  Barwise's route passes through
      the Karp property (\cref{thm:p03-w03-p03-048-lowenheim-karp}), the
      boundedness of \(\mathcal L^{\ast}\) by \(\omega\)
      (\cref{thm:p03-w03-p03-048-karp-omega}) and the \(\Sigma^1_1\)
      definability of \((\omega,<)\) (\cref{thm:p03-w03-p03-048-finite-omega});
      the last step uses an indexed union, which the host package does not
      provide, and \cref{thm:p03-w03-p03-048-host-route} replaces it by a
      direct use of compactness.  Barwise's relation \(\cong_{p}\), for partial isomorphisms
      of any size, and the host's potential isomorphism, for finite partial
      isomorphisms, agree between countable relational structures, since
      both are then equivalent to isomorphism
      (\cref{thm:p03-w03-p03-048-countable}).
\item \emph{Conclusions.}  Both conclude that every sentence is equivalent,
      over its occurrence language, to a first-order sentence in its
      symbols.  Barwise's conclusion covers the empty structure and
      function symbols.
\end{enumerate}
\end{hbconvention}


The next two results run Barwise's Karp route inside the regular relational
package itself.  There a partial isomorphism between \(\tau\)-structures is a
one-one map between subsets of the universes that preserves and reflects
every relation of \(\tau\) on its domain; this is clause~3 of
\cref{def:p03-w03-p03-048-standing} for a language without function
symbols.  A nonempty set \(I\) of partial isomorphisms with the back and
forth property is written \(I:A\cong_{p}B\) as before.  Potential
isomorphism (\cref{def:p03-potential-isomorphism}) is the case in which the
members of \(I\) are finite.

\begin{hblemma}{thm:p03-w03-p03-048-host-karp}{proved=annex:p03-w03-p03-048-host-karp}{Barwise's Theorem~2.2 in the regular relational package}
Let \(\mathcal L^{\ast}\) be a regular abstract logic in the sense of
\cref{sec:p03-lindstrom} satisfying hypothesis~2 of
\cref{thm:p03-lindstrom}.  If \(A\) and \(B\) are \(\tau\)-structures and
\(I:A\cong_{p}B\), then \(A\) and \(B\) satisfy the same
\(\mathcal L^{\ast}\)-sentences of \(\tau\).  In particular potentially
isomorphic structures do.
\end{hblemma}

\begin{hbtheorem}{thm:p03-w03-p03-048-host-route}{proved=annex:p03-w03-p03-048-host-route}{\Cref{thm:p03-lindstrom} along the Karp route}
Let \(\mathcal L^{\ast}\) be a regular abstract logic satisfying
hypotheses~1 and 2 of \cref{thm:p03-lindstrom}, and let \(\varphi\) be an
\(\mathcal L^{\ast}\)-sentence with finite occurrence signature
\(\tau_\varphi\).
\begin{enumerate}
\item (Barwise's 3.2, \((\Leftarrow)\), construction.)  If \(\varphi\) is
      equivalent over \(\tau_\varphi\) to no relational first-order sentence,
      there are a finite relational signature \(\sigma\) with a unary \(U\)
      and a binary \(<_U\), and an \(\mathcal L^{\ast}\)-sentence \(\psi\) of
      \(\sigma\), such that in every model of \(\psi\) the relation \(<_U\)
      well-orders \(U\), and for every \(m\geq1\) some model has
      \((U,<_U)\) of order type \(m\).
\item (Compactness in place of Barwise's 1.4, 1.2 and 3.4.)  No
      \(\mathcal L^{\ast}\)-sentence has the two properties in part~1.
\end{enumerate}
Consequently every \(\mathcal L^{\ast}\)-sentence is equivalent over its
occurrence signature to a relational first-order sentence, which is the
conclusion of \cref{thm:p03-lindstrom}.
\end{hbtheorem}

The proof of part~1 uses hypothesis~2 only through
\cref{thm:p03-w03-p03-048-host-karp}, and part~2 uses hypothesis~1 only.
Barwise's own route continues from part~1 through his Theorem~1.4 (the
supremum \(\omega\) of the finite types is accessible) and his Theorem~1.2
(\((\omega,<)\) is \(\Sigma^1_1\)), and only then applies countable
compactness.  Both of those steps form indexed unions by a \(1\)-term
interpretation, which the relational package does not provide; part~2
instead applies compactness to \(\psi\) directly, in the way that the
boundedness half of \cref{thm:p03-w03-p03-048-karp-omega} applies
first-order compactness.  The upward L\"owenheim--Skolem alternative~(a) of
Barwise's Corollary~3.4 is not carried over: its proof uses the
\(\Sigma^1_1\) definition of \((\omega,<)\), and so the indexed unions.
\begin{hbexample}{def:p03-w03-p03-048-example}{The sentence of Theorem~2.2 for \((\mathbb Q,<)\) and \((\mathbb R,<)\)}
Let \(L=\{\mathbf V,<\}\), \(\mathfrak M_0=(\mathbb Q,<)\),
\(\mathfrak M_1=(\mathbb R,<)\), and let \(I\) be the set of finite
order-preserving maps from \(\mathbb Q\) to \(\mathbb R\).  \(I\) is
nonempty and has the back and forth property, because between and beyond
finitely many reals there are rationals and conversely; so
\(I:\mathfrak M_0\cong_{p}\mathfrak M_1\).  In the structure built in the
proof of \cref{thm:p03-w03-p03-048-lowenheim-karp}, \(U=\mathbb Q\),
\(W=\mathbb R\), the positions are the members of \(I\), and \(E(s,x,y)\)
says that the map \(s\) sends \(x\) to \(y\).  If some sentence \(\varphi\)
of a L\"owenheim logic were true in \((\mathbb Q,<)\) and false in
\((\mathbb R,<)\), the countable model of the combined sentence would
contain two countable dense orders without endpoints, one satisfying
\(\varphi\) and one not, joined by a back and forth system; they are
isomorphic by \cref{thm:p03-w03-p03-048-countable}, against the Isomorphism
Axiom.  So no L\"owenheim logic in Barwise's sense separates
\((\mathbb Q,<)\) from \((\mathbb R,<)\); the same pair shows in
\cref{def:p03-w03-p03-039-q-r} that countability cannot be dropped from the
back and forth lemma.
\end{hbexample}

\begin{hbwarning}{def:p03-w03-p03-048-limits}{What the route does not give}
The Karp property by itself does not make a logic first-order:
\(\mathcal L_{\infty\omega}\) has it \cite[Part~I, \S10, 10.3, p.~244;
Part~II, p.~256]{Barwise1974Axioms}, and Theorem~3.2 locates the missing
ingredient in the boundedness of \(\mathcal L^{\ast}\) by \(\omega\).  By
the arguments for \cref{thm:p03-w03-p03-048-lindstrom}, a logic in which
\((\omega,<)\) is \(\Sigma^1_1\) fails (b), and fails (a) as well if it has
the L\"owenheim property.  Barwise reports that \((\omega,<)\) is definable
in each logic of his Part~I, \S1, except \(\mathcal L_{\omega\omega}\) and
\(\mathcal L(Q_1)\) \cite[Part~II, \S1, p.~250]{Barwise1974Axioms}; for the
quantifier ``there exist infinitely many'' compare
\cref{thm:p03-w03-p03-047-q0-omega}.  The proof of alternative~(a) needs the L\"owenheim property
twice, once for the Karp property and once to obtain a countable model of
the \(\Sigma^1_1\) definition of \((\omega,<)\).
\end{hbwarning}
